
This study tested whether annotators with prior exposure to deployed RLHF-trained AI systems prefer AI-idiomatic stylistic features more strongly than annotators without such exposure, using the Anthropic HH-RLHF dataset and a conditional logit model with semantic cluster fixed effects.

The AIFS construct — a composite count across four regex-detected feature classes — was found to be a valid preference predictor: β₁ = +0.0246 (p < 0.001) across 80,342 preference pairs and 27,034 semantic clusters, with controls for response length and cluster heterogeneity. This establishes AIFS as an operational measure for use in future preference modeling studies.

The adaptation hypothesis was not confirmed. The interaction coefficient β₄ = −0.0016 (OR = 0.9984, 95\% CI [0.9861, 1.0108], p = 0.796) is not statistically significant and does not satisfy any of the four pre-registered gate thresholds. The 95\% CI width of 0.025 is narrow enough to exclude practically meaningful effects, indicating that the null is not attributable to insufficient power given the proxy variable used.

Examination of the dataset reveals that the null result cannot be interpreted as evidence for or against the adaptation hypothesis because the proxy variable — the helpful-base versus helpful-online split — encodes response generation provenance rather than annotator exposure history. The split is a valid data provenance label but not a valid annotator condition variable.

The primary implication is methodological: empirical testing of the bidirectional alignment hypothesis requires preference corpora that include annotator-level metadata. Existing public RLHF datasets, including HH-RLHF, do not provide this. The AIFS pipeline and the conditional logit infrastructure developed here are directly reusable when appropriate annotator-linked datasets become available.

